\section{Related Work}
\label{sec:related}

\para{Representation engineering.}
\citet{zou2023representation} introduced Linear Artificial Tomography (LAT), showing that high-level concepts---honesty, harmfulness, emotions---are linearly encoded in LLM hidden states and can be extracted via PCA on contrastive activation differences.
\citet{marks2024geometry} extended this to truth/falsehood, demonstrating that simple difference-in-means probes generalize across topically diverse datasets and are more causally implicated than complex classifiers.
\citet{arditi2024refusal} showed that refusal behavior is mediated by a single direction in activation space, and that ablating this direction eliminates refusal with minimal capability degradation.
These results establish that safety-relevant behaviors can be concentrated in low-dimensional subspaces---a finding we extend to the novel domain of offense.

\para{Emotion and affect in LLM representations.}
Several recent studies have probed the emotional landscape of LLMs.
\citet{dipalma2025llamas} found that binary sentiment is best detected in mid-layers ($\sim$60--75\% depth) with 94--96\% probe accuracy, while fine-grained 6-class emotion classification is harder and peaks in earlier layers.
\citet{tak2025mechanistic} provided causal evidence for emotion processing in mid-layer attention heads, showing that the appraisal--emotion structure in LLMs mirrors cognitive appraisal theory from psychology.
\citet{reichman2025emotions} discovered a universal emotional manifold via SVD, with principal components corresponding to valence, dominance, approach--avoidance, and arousal.
Most directly relevant to our work, \citet{keeman2026whether} identified a fundamental dissociation between \emph{affect reception} (detecting that content is emotionally significant, AUROC = 1.000 even on keyword-free stimuli) and \emph{emotion categorization} (identifying which emotion, partially keyword-dependent).
This dissociation predicts that a binary offense probe should achieve near-perfect accuracy, while fine-grained offense severity estimation will be substantially harder---a prediction our results confirm.

\para{Probe reliability and Goodhart's Law.}
A growing body of work raises concerns about the reliability of representation probes under optimization pressure.
\citet{bailey2025obfuscated} showed that all evaluated latent-space defenses---probes, sparse autoencoders, out-of-distribution detectors---are vulnerable to adversarial obfuscation attacks that reduce probe recall from 100\% to 0\% while retaining jailbreak success.
\citet{wehner2025probe} found that when probes are used as training signals, DPO-based preference learning preserves probe detectability, while SFT-based gradient regularization allows partial evasion---an unexpected asymmetry with implications for probe-based safety systems.
\citet{braun2025understanding} demonstrated that steering vectors exhibit high variance across samples and can produce effects opposite to the desired direction, particularly when the target behavior lacks coherent representation.
These findings motivate our focus on convergent validity: rather than relying on a single probe, we triangulate across multiple methods and behavioral evidence to assess trustworthiness.

\para{Positioning.}
Unlike prior work that studies toxicity (a property of content) or emotions (a general affective state), we study \emph{taking offense}---a second-person affective reaction.
Unlike probe reliability studies that focus on adversarial robustness \citep{bailey2025obfuscated} or Goodhart dynamics \citep{wehner2025probe}, we address the more fundamental question of what happens when \emph{no ground truth exists} for the target concept.
Building on the affect reception framework of \citet{keeman2026whether}, we test whether the binary/fine-grained dissociation extends to offense and use it to establish bounds on probe trustworthiness.